\section{Comparison with existing labelling designs}\label{sec:baselines}

PATTERN-DIF combines four ingredients that are each known: sampling where models disagree, stratification on a model-based proxy with Neyman-type allocation, design-based ratio estimation and a model-assisted (difference) estimator. This section asks whether natural alternatives built from one or two of these ingredients already do as well. The study was run before the final form of the method was fixed. Its PATTERN-DIF variant (B5, $\lambda=1$) is the same estimator as the main method, but the samples differ from those of Section~\ref{sec:empirical} (different random seeds). Values differ between the two studies by a median of about 5\% per cell (up to 27\% in single cells, law school excluded) and the medians over settings by up to 14\% (PATTERN-DIF at $L=250$: 6.6 here, 5.8 in Section~\ref{sec:empirical}). The abstract and the main results use the study of Section~\ref{sec:empirical}. We ran the baselines and SRS and PATTERN on the eleven settings of Section~\ref{sec:empirical}, at $L\in\{250,500,1000,2000\}$ with $L\le N/4$ (35 setting--budget cells) and $R=200$ replications per cell. To isolate the effect of labelling and estimation, every method estimates the same totals $T$, maps them through the same $\hat\theta=f(\hat T)$ to the same 108 comparison estimands, uses the same proxy $\pi_i$ and allocation weights where it has any, and reports 95\% Wald intervals; the budget is counted in distinct labelled instances. Estimands that every design estimates exactly are excluded (51 of 108 in Law school, none elsewhere). Full definitions are in Supplementary Section~\SIref{si:baselines}.

\paragraph{Baselines.} B1 is importance sampling for active comparison \citep{sawade2012active}: draws with replacement from an instrumental distribution that is proportional to the root of the proxy-expected squared influence, adapted here to TPR and FPR gaps, with self-normalised totals. B2 stratifies on the proxy score (group $\times$ quantile bins of $\pi_i$; Score-10 and Score-5) with the same proxy-Neyman allocation as PATTERN \citep{bennett2010online,kumar2018classifier}, so it differs from PATTERN only in the strata. B3 is a disagreement-focused two-phase design (Dis-25, Dis-50): a simple random first phase of $L/4$ or $L/2$ labels, then a simple random sample of the disagreement set for the rest \citep{sawade2012active,musgrove2023discordant}. B4 is prediction-powered inference: PPI with $\lambda=1$ \citep{angelopoulos2023ppi}, PPI++ with a tuned $\lambda$ \citep{angelopoulos2026ppipp} and StratPPI, which applies PPI++ within proxy strata with a Neyman-type allocation \citep{fisch2024stratppi}. B5 is the proposed method: the PATTERN design with the stratified difference estimator ($\lambda=1$, PATTERN-DIF), and a tuned-weight variant (PATTERN-PPI).

\begin{table}[t]\centering\small
\caption{Baselines by label budget $L$ (number of settings in brackets). Efficiency: ratio of mean squared errors, median over the non-degenerate comparison estimands within a setting, then median over settings; ``vs SRS'' is MSE(SRS)/MSE(method), ``vs PATTERN'' is MSE(PATTERN)/MSE(method), so values above 1 favour the method. Coverage: nominal 95\% Wald intervals, mean over estimands and settings; ``worst'': lowest single-estimand coverage over all settings and budgets. $R=200$ replications per setting and budget. Source: \texttt{results/label\_baselines.json}.}\label{tab:base-budget}
\setlength{\tabcolsep}{3.5pt}
\resizebox{\textwidth}{!}{\begin{tabular}{lrrrrrrrrrrr}\toprule
 & \multicolumn{2}{c}{$L=250$ (11)} & \multicolumn{2}{c}{$L=500$ (10)} & \multicolumn{2}{c}{$L=1000$ (8)} & \multicolumn{2}{c}{$L=2000$ (6)} & \multicolumn{3}{c}{All $(s,L)$} \\
\cmidrule(lr){2-3} \cmidrule(lr){4-5} \cmidrule(lr){6-7} \cmidrule(lr){8-9} \cmidrule(lr){10-12}
Method & vs SRS & vs PAT. & vs SRS & vs PAT. & vs SRS & vs PAT. & vs SRS & vs PAT. & Cover. & Worst & ms \\ \midrule
SRS & 1 & 0.27 & 1 & 0.25 & 1 & 0.23 & 1 & 0.14 & 0.925 & 0.16 & 8 \\
PATTERN & 3.75 & 1 & 3.95 & 1 & 4.35 & 1 & 7.43 & 1 & 0.947 & 0.52 & 8 \\
IS (B1) & 3.09 & 0.79 & 2.77 & 0.74 & 3.48 & 0.67 & 4.30 & 0.54 & 0.950 & 0.60 & 9 \\
Score-10 (B2) & 1.45 & 0.40 & 2.09 & 0.52 & 2.24 & 0.56 & 4.16 & 0.50 & 0.941 & 0.17 & 8 \\
Score-5 (B2) & 1.97 & 0.54 & 2.24 & 0.57 & 2.15 & 0.52 & 3.18 & 0.43 & 0.943 & 0.37 & 8 \\
Dis-25 (B3) & 2.25 & 0.42 & 2.27 & 0.55 & 2.04 & 0.35 & 2.48 & 0.26 & 0.942 & 0.45 & 9 \\
Dis-50 (B3) & 1.84 & 0.43 & 1.94 & 0.37 & 1.68 & 0.34 & 1.93 & 0.25 & 0.940 & 0.36 & 9 \\
PPI (B4) & 1.82 & 0.49 & 1.82 & 0.47 & 1.86 & 0.44 & 1.81 & 0.27 & 0.940 & 0.27 & 34 \\
PPI++ (B4) & 2.10 & 0.49 & 1.97 & 0.48 & 1.93 & 0.44 & 1.89 & 0.27 & 0.929 & 0.16 & 49 \\
StratPPI (B4) & 3.07 & 1.01 & 4.45 & 1.12 & 4.58 & 1.08 & 8.50 & 0.83 & 0.924 & 0.14 & 44 \\
\midrule \multicolumn{12}{l}{\emph{Extension (B5): PATTERN sample, model-assisted estimator}} \\
PATTERN-DIF (B5) & 6.57 & 1.63 & 6.23 & 1.48 & 6.64 & 1.48 & 15.5 & 1.32 & 0.946 & 0.64 & 35 \\
PATTERN-PPI (B5) & 6.71 & 1.68 & 6.37 & 1.55 & 6.83 & 1.54 & 15.5 & 1.33 & 0.929 & 0.20 & 51 \\
\bottomrule\end{tabular}}\end{table}


\paragraph{Results.} Table~\ref{tab:base-budget} summarises all methods by budget. PATTERN is the most efficient of the designs that use the same plain estimator as SRS (B1 to B3) at every budget: its median efficiency against SRS is 3.75, 3.95, 4.35 and 7.43 at $L=250,500,1000,2000$. It is not the most efficient method overall: StratPPI, which combines proxy strata with a model-assisted estimator, matches or exceeds it in the paired medians at $L=250$, 500 and 1000 (ratios 1.01, 1.12, 1.08) and has the higher median efficiency against SRS at $L=500$, 1000 and 2000 (4.45, 4.58, 8.50), and the extensions B5, which add the same kind of correction to the PATTERN sample, beat both. Among the competing designs the closest are importance sampling (B1: 3.09, 2.77, 3.48, 4.30 against SRS; 0.54--0.79 of PATTERN) and StratPPI (B4: 3.07, 4.45, 4.58, 8.50; 1.01, 1.12, 1.08, 0.83 of PATTERN). Proxy-score strata alone (B2) reach 1.45--4.16 against SRS, the disagreement-only designs (B3) 1.7--2.5 and PPI/PPI++ with simple random labels 1.8--2.1, all clearly below PATTERN (medians 0.25--0.57 of its efficiency). PPI and PPI++ with SRS labels barely move with $L$: they use the proxy but not the fact that the numerators vanish where models agree. The model-assisted extensions improve on PATTERN in every cell except Law school at $L=2000$, where PATTERN is already exact, by a median factor of 1.3--1.7 depending on budget and variant and up to 11 in Law school at $L\le500$; PATTERN-DIF reaches 6.6, 6.2, 6.6 and 15.5 against SRS. The estimated-$\lambda$ version adds little (within 3\% of PATTERN-DIF in the median) and costs coverage, so we recommend the fixed-$\lambda$ version. The extension changes the estimator and leaves the design alone. Its design-unbiasedness and variance estimator are covered by the theory of Section~\ref{sec:theory}, which treats the difference estimator with a fixed coefficient; the tuned-coefficient version is covered only by the asymptotic equivalence statement there.

\paragraph{Where a baseline beats PATTERN.} Supplementary Table~\SIref{tab:base-setting} and Supplementary Figure~\SIref{fig:base-efficiency} (Supplementary Section~\SIref{si:baselines}) show per-setting ratios. First, StratPPI beats PATTERN at $L=500$ in seven of the ten settings of the table (ratios 1.01--1.57) and is below it in Adult-sex (0.97), Adult-race (0.93) and Churn (0.24); over all cells its ratio exceeds 1 in 19 of 35. The margin is small in most settings and reflects the model-assisted estimator that StratPPI uses: the same correction applied to the PATTERN sample improves PATTERN itself by a median factor of 1.3--1.7. Second, Dis-25 beats PATTERN on Churn at $L=500$ (ratio 1.52), where the disagreement set is only 19\% of the pool, but its intervals are unreliable there (mean coverage 0.93, worst-estimand coverage 0.66). Third, PPI++ beats PATTERN on EEG (ratio 1.19 at $L=500$), the setting with the largest disagreement set (73\%). Importance sampling never beats PATTERN in the all-estimand median (ratio below 1 in all 35 cells), but on the pairs among the three most accurate models it is competitive: its median efficiency against SRS on those pairs is 4.2, 4.1, 4.2 and 6.3 against 3.4, 3.6, 4.0 and 6.8 for PATTERN, and it wins on those pairs in 10 of 35 cells. If only the best few models matter, importance sampling is a serious alternative. Conversely, in the two settings where PATTERN gains most (Law school and Churn, whose disagreement sets cover only 19\% of the pool) it is 6--84 times more efficient than SRS at $L\le1000$ and labels essentially every disagreement instance at $L=2000$ in Law school, because it gives the agreement region a stratum of its own that can be labelled sparsely.

\paragraph{Coverage.} Over all 35 cells the mean 95\% coverage is 0.925 (SRS), 0.947 (PATTERN), 0.950 (IS), 0.941 (Score-10), 0.942 (Dis-25), 0.929 (PPI++), 0.924 (StratPPI), 0.946 (PATTERN-DIF) and 0.929 (PATTERN-PPI). Shortfalls concentrate in the smallest cells. The lowest worst-estimand coverages are 0.14 for StratPPI (Churn, $L=250$), 0.16 for SRS and PPI++, 0.20 for PATTERN-PPI and 0.52 for PATTERN. The tuned estimators undercover when per-stratum weights are estimated from a few labels; PATTERN-DIF, which fixes $\lambda=1$, keeps the coverage of PATTERN. Undefined estimates occur in at most 1\% of replications for any method. Details are in Supplementary Section~\SIref{si:baselines}.

\paragraph{Limitations.} The comparison is restricted to fixed-budget, one-shot designs on cached pools of nine models. The proxy is the same for every method, so a better proxy could change the ranking, especially for StratPPI and PATTERN-PPI, which rely on it directly. The importance-sampling baseline draws with replacement and so spends a median of $1.12L$ draws for $L$ labels; counting draws instead of labels would lower its efficiency. With $R=200$ replications the Monte Carlo error of a single mean squared error is about 10\%, so ratios between 0.9 and 1.1 should not be read as a ranking, and medians over estimands hide differences between estimands.
